% Complete proof independently audited by root and agents B and D.
\subsection{Small act cones under finite-target nullity}

\begin{lemma}[The precise finite-target uses in the small band hulls]
\label{C:small-band-finite-target-cups}
Let $P$ be a finite monoid with rectangular-band minimal ideal.
Assume that every integral class in $H^2(BP;\mathbb Z)$ restricts
to zero along every map from the classifying space of a finite group.
Then the following consequences of the earlier hull calculations hold.
\begin{enumerate}
\item For a $3\times3$ band, all rational degree-two products vanish.
\item For a $3\times4$ band with rational second-cohomology rank two,
all rational degree-two products vanish.
\item For a $3\times4$ band with rational second-cohomology rank one,
the square of an integral generator cannot be nonzero and primitive modulo torsion
in fourth cohomology.
\end{enumerate}
The same statements hold after opposition. There is no finite-CW-type
assumption on $BP$ in this lemma.
\end{lemma}
\begin{proof}
Read the proofs of Theorems~\ref{B:three-by-three-cups} and
\ref{B:three-by-four-targets} with the stated nullity hypothesis.
Every use there of finite CW type is to rule out an actor with a
nonzero integral character on its eventual cyclic group. A preimage
actor has an eventual finite group surjecting onto that cyclic group;
inflation of the character is injective. The corresponding integral
hull class pulls back to a class on $BP$, contradicting the present
hypothesis. All other cases use actual cup primitives, rank-one
square primitives, or zero fourth rational homology in a containing
hull. Those conclusions pull back without any finiteness premise.

The only nonzero-square case for a $3\times4$ band is the rank-one
even-Hopf hull. Its integral second homology is $\mathbb Z$, and the
square of its primitive degree-two class evaluates evenly on every
integral four-cycle. The balanced-tensor presentation makes the map
from $P$ onto this hull surjective on integral second homology.
Even if $H_2(BP;\mathbb Z)$ has torsion, its free rank is one,
and a surjection to $\mathbb Z$ is an isomorphism on the free part.
The hull generator therefore pulls back to a generator on $BP$.
Its square still evaluates evenly, which prevents primitivity modulo
torsion. In rank two this exceptional hull cannot occur at all.
This proves the three assertions, including the integral precision.
\end{proof}

\begin{lemma}[Nullity survives the projective orbit quotient]
\label{C:orbit-quotient-local-nullity}
Let $X=B\mathcal C(M,S)$ be simply connected, with minimal band
$I\times\{0,1\}$ and no band-fixed state. Suppose every integral
second-cohomology class of $X$ restricts to zero on every finite $BG$.
The constant-action quotient and its compression $Y\simeq BP$ have
the same nullity property in degree two.
\end{lemma}
\begin{proof}
Corollary~\ref{C:two-column-orbit-wedge} gives
$Y\simeq X\vee\bigvee_{o-1}S^2$, under the actual quotient map.
Every degree-two class on this wedge is a sum of pullbacks from the
two projections to $X$ and to $\bigvee_{o-1}S^2$.
The first restrictions vanish by hypothesis. The second vanish by
Miller's theorem with the finite target $\bigvee_{o-1}S^2$.
Thus every such restriction vanishes on $Y$ as asserted.
\end{proof}

\begin{lemma}[The rational five-outside row argument at six states]
\label{C:rational-four-row-five-outside-six-states}
Let $M$ have minimal band $4\times2$ and exactly five outside
elements. An act cone on five or six states cannot have the rational
homology and cohomology ring of $\mathbb CP^2$.
\end{lemma}
\begin{proof}
The fixed-state case has zero degree-two products. Otherwise the
projective orbit kernel gives a constant-action compression with the
same rational homology in degrees at least three. If $o$ is the
number of core orbits, its row and column modules satisfy
\[
 \dim V=3,\qquad \dim W=|S|-2o+1,
 \qquad \dim\operatorname{Tor}_{0,1,2}^A(W,V)=(o,0,1),
\]
and higher Tor vanishes. Here $A$ is the five-dimensional contracted
algebra of $M$. This uses only the rational homology assumptions.

The local case is the universal affine obstruction
Lemma~\ref{C:local-two-column-act-cups}. The nonlocal classification
and the complete module tables used in
Theorem~\ref{C:four-row-five-outside-seven-states} leave the line
and Nakayama algebras. For Nakayama, the only possible row module
is $S_2\oplus Ae_2$, and its sole quotient character occurs twice;
the missing-character or two-covector affine-origin lemma gives zero
products. For the line, the split-character row and the row with a
twice-occurring quotient character are excluded by the same universal
arguments. The remaining row is $V=S_3\oplus L_a$, and the module
table gives
\[
 W=T_1\oplus T_3^{\oplus t}\oplus R_a^{\oplus p}
                         \oplus R_b^{\oplus q},\qquad
 o=t+p+q,\qquad |S|=4o-t,
 \qquad 0\leq t\leq o.
\]
Five states admit no solution: $o=1$ gives at most four states,
whereas $o\geq2$ gives at least six. For six states the only
possibility is $o=2,t=2,p=q=0$. Both balanced degree-two forms
then factor through the single $S_3$ row summand and have matrix
rank one. Their squares are zero by Lemma~\ref{B:rank-one-square}.
Their pullbacks span the one-dimensional second cohomology of the
original cone; one pulls back nontrivially and contradicts its
projective-plane square. No finite model for the compression is used.
\end{proof}

\begin{theorem}[A finite-target obstruction through twenty act arrows]
\label{C:finite-target-small-act-twenty}
Let $X=B\mathcal C(M,S)$ for a nonempty finite right act, with
$|M|+|S|+1\leq20$. The following conditions cannot hold simultaneously:
\begin{enumerate}
\item $X$ is simply connected and its rational homology is that of
$\mathbb CP^2$;
\item an integral generator $v\in H^2(X;\mathbb Z)\cong\mathbb Z$
has square nonzero and primitive modulo torsion in $H^4(X;\mathbb Z)$;
\item for every finite group $G$ and every map $BG\to X$, the
restriction of $v$ is zero;
\item over $\mathbb F_{65521}$, the Betti numbers in degrees
two, three and four are $(1,0,1)$.
\end{enumerate}
In particular no finite-CW-type assumption on $X$ is required.
\end{theorem}
\begin{proof}
Write the minimal Rees ideal as $I\times H\times J$, put
$r=|I|$, $c=|J|$, $q=|M|-rc|H|$, and $s=|S|$.
The nonzero square requires $H^2(BM;\mathbb Q)\ne0$; otherwise
every degree-two class is relative to $BM$ and its products vanish
in the suspension quotient. The rational balanced presentation
therefore gives $r,c\geq2$. The all-order two-row act theorem gives
$r\geq3$. The finite-line argument gives $s\geq3$. The identity
lies outside the minimal ideal in these cases, so $q\geq1$.

If $H$ is nontrivial, the arrow budget leaves only
$r=3,c=2,|H|=2,q\leq4$. Lemma~\ref{C:three-row-forced-profiles},
which permits a nontrivial minimal group, requires at least five
outside elements for a nonzero product. Thus $H$ is trivial below.
The original square is nonzero rationally by its primitivity.

For every compression used below, an actual homotopy equivalence
preserves the generator and primitive square. If an orbit quotient
with two-column projective kernel is used instead, its higher
rational homology is unchanged, its second-cohomology rank is the
number $o$ of core orbits, and
Lemma~\ref{C:orbit-quotient-local-nullity} supplies nullity for all
its integral degree-two classes. We explicitly indicate these uses.

\emph{Three states.}
If a band actor is nonconstant, the constant-action quotient has at
most two states. Its degree-two cohomology surjects onto that of the
original cone, and its finite-line products are zero. Otherwise
compression gives an actual equivalent monoid with band $r\times3$
and $q$ outside elements. Here $|M|\leq16$.
For $r=3$ use the first part of
Lemma~\ref{C:small-band-finite-target-cups}; for $r=4$ use its
primitive-square part. For $r=5$, its order is
$15+q\leq15+(16-2r)=21$, so
Proposition~\ref{B:three-by-five-finite-target-transfer} applies after
opposition. If $r\geq6$, the same inequality gives $q\leq4$,
and Theorem~\ref{B:contracted-four-Tor} contradicts bounded
rational homology with nonzero fourth homology.

\emph{Four states.}
Lemma~\ref{C:four-state-band-dichotomy} uses only first rational
homology vanishing. It gives zero products or constant band action.
In the latter case the equivalent compression has band $r\times4$,
with $|M|\leq15$. Three rows are excluded by the primitive-square
part of Lemma~\ref{C:small-band-finite-target-cups}.
For four rows and two old columns, the compressed order is
$16+q\leq23$; apply
Proposition~\ref{B:four-by-four-finite-target-transfer}.
For four rows and at least three old columns, $q\leq3$, so the
bounded Tor obstruction already applies. For five rows, necessarily
$c=2$ and $q\leq5$. Values $q\leq4$ again give that obstruction;
at $q=5$ the compressed $5\times4$ band has augmentation dimensions
four and three. The proof and complete module table of
Lemma~\ref{C:four-by-five-five-outside-CP2} exclude the rational
Tor sequence $(1,0,1,0,\ldots)$ without any finiteness premise.
For at least six rows, $q\leq3$ and
Lemma~\ref{C:three-outside-all-two-column-act-cups} applies.

\emph{At least five states and two columns.}
For three rows the budget gives $q\leq8$, hence $|M|\leq14$.
The bounded-rational assertion of
Theorem~\ref{C:three-row-fourteen-act-cones} excludes the square.
For at least five rows, $q\leq4$ and the stronger rational
assertion of Theorem~\ref{C:two-column-four-outside-act-cones}
applies at every state count.
It remains that $r=4$, where $q+s\leq11$.
The cases $q\leq4$ have the same Tor obstruction; $q=5$ gives
$s=5$ or $6$ and the preceding rational module lemma applies.
The sole remaining possibility is $q=6,s=5$.

In this last case a fixed state makes the products zero.
If the band action is nonconstant, the connected act has exactly
two two-point core orbits and one additional state. Its quotient
has three states. The projective orbit kernel and actual wedge
equivalence give a compression with band $4\times3$, rational
second-cohomology rank two, and nullity of all integral degree-two
group restrictions. The rank-two clause of
Lemma~\ref{C:small-band-finite-target-cups} makes its products
zero, and pullback gives the contradiction.

If instead the band acts constantly, use the exact source census
and row-module partition in
Lemma~\ref{C:four-row-six-outside-source-census} and
Proposition~\ref{C:four-row-six-outside-row-reduction}.
Projective-dimension-one rows give zero products universally;
the stable and local-corner rows give zero fourth rational homology
under boundedness. For the remaining $173$ rows, the independently
verified exact projective resolutions over $\mathbb F_{65521}$ and
the exhaustive five-state act calculation in
Theorem~\ref{C:four-row-six-outside-six-state-cones} apply.
All $110418$ constant five-state acts fail the necessary tensor
dimensions $(1,0,1)$ in Tor degrees zero, one and two.
These are the actual Betti numbers in degrees two, three and four
of the compressed cone. Condition (4) is therefore impossible.
This uses exact field resolutions and a field homology obstruction,
not an inference from a rank measured at an uncontrolled prime.

\emph{At least five states and at least three columns.}
The budget leaves only bands $3\times3$, $3\times4$, and $4\times3$.
For either twelve-element band, $q+s\leq7$. Thus $q=1$ gives a
projective row module; otherwise $q=2$ and $s=5$. If the action is
constant, bounded rational Tor with $q\leq4$ is impossible.
If it is not, connectedness and the absence of a fixed state give
at least two core orbits, at least one of which has two points.
Its column map makes the old-column image in the constant-action
quotient have at most two points. The universal transfer in
Lemma~\ref{C:old-column-image-two-transfer}, followed by
Lemma~\ref{C:three-outside-all-two-column-act-cups}, kills products.

For the $3\times3$ band, $q+s\leq10$.
An old-column image of size one gives zero old/old affine cochains.
Size two gives a literal common compression over a $3\times2$
monoid with the same $q\leq5$ outside elements, whose act products
vanish universally by Theorem~\ref{C:three-row-fourteen-act-cones}.
If all three old-column values remain distinct, every core orbit
has three points. Lemma~\ref{C:full-column-orbit-kernel} then
preserves bounded higher rational homology under the quotient.
For $q\leq4$ use Theorem~\ref{B:contracted-four-Tor}.
The only other possibility is $q=5,s=5$. Two full core orbits
cannot fit, so the band already acts constantly. Compression is an
actual equivalence to a $3\times5$ band monoid of order twenty,
excluded by Proposition~\ref{B:three-by-five-finite-target-transfer}.
This exhausts the arrow budget and proves the theorem.
\end{proof}

\begin{corollary}[Nontrivial projective opposite corners through twenty-one arrows]
\label{C:projective-opposite-corner-twenty-one}
Let $C$ be a triangular two-object category with endomorphism monoids
$M,N$, cross-biset ${}_NS_M$, and projective trivial right
$\mathbb QN$-module. If $|N|\geq2$ and $|C|\leq21$, then
$BC\not\simeq\mathbb CP^2$.
The same conclusion holds with $|C|\leq22$ and $|N|\geq3$.
Both statements also hold after opposition.
\end{corollary}
\begin{proof}
Assume otherwise. The actual orbit quotient of
Theorem~\ref{C:finite-target-opposite-corner} has at most
$|C|-|N|+1\leq20$ arrows.
Corollary~\ref{C:finite-target-corner-CP2} supplies simple
connectivity, rational projective-plane homology, the primitive
integral square, and nullity of every integral degree-two class
on every finite group. The group order of the opposite corner
is at most $22$, so $65521$ is a good prime and the same comparison
supplies condition (4). The preceding theorem gives a contradiction.
\end{proof}

\begin{remark}[The remaining two-object scope]
\label{C:finite-target-small-act-scope}
The argument treats the precise finite-target hypotheses above,
rather than declaring every rational projective-plane act cone to
have finite CW type. At twenty-one arrows, an unexcluded triangular
category must fail both projective-corner tests: its source corner
must have at least two minimal rows and its target corner at least
two minimal columns. This follows by applying the corollary also
to the opposite category. It does not dispose of these nonprojective
opposite corners. Strongly connected two-object categories are resolved
separately by Theorem~\ref{C:twentytwo-strongly-connected}; the
present theorem itself concerns only triangular categories. The
original unrestricted category minimum is not asserted here.
\end{remark}
